% =============================================================================
%  Tài liệu 01 - Hồ sơ khởi động dự án (Project Initiation Document)
%  Build: latexmk 01-project-initiation.tex   (pdfLaTeX, output in build/)
% =============================================================================
\documentclass[12pt,a4paper]{report}
\usepackage{sumtalk-report}
% =============================================================================
%  sumtalk-info.tex - information shared by every document (01-06).
%  Edit here once; every cover and footer picks it up.
% =============================================================================
\newcommand{\ProjectName}{SumTalk}
\newcommand{\CourseName}{QUẢN LÝ DỰ ÁN PHẦN MỀM}          % cover (upper case)
\newcommand{\CourseNameNormal}{Quản lý Dự án Phần mềm}     % footer
\newcommand{\AcademicYear}{2026 -- 2027}
\newcommand{\Instructor}{\canDien[họ tên giảng viên]}

% One \item per member: "Full name -- student ID", sorted by given name.
\newcommand{\TeamMembers}{%
  \item Nguyễn Anh Duy -- \canDien[MSSV]
  \item Đỗ Thành Đạt -- \canDien[MSSV]
  \item Trần Hải Đăng -- \canDien[MSSV]
  \item Châu Hùng Hào -- \canDien[MSSV]
  \item Bùi Đoàn Minh Nhật -- \canDien[MSSV]
  \item Võ Hoàng Sơn -- \canDien[MSSV]
}


\newcommand{\DocTitle}{HỒ SƠ KHỞI ĐỘNG DỰ ÁN}
\newcommand{\DocTitleNormal}{Hồ sơ khởi động dự án}
\newcommand{\DocTitleEn}{Project Initiation Document}
\newcommand{\DocRunningTitle}{Hồ sơ khởi động dự án}
\newcommand{\DocCode}{SUMTALK-PID}
\newcommand{\DocVersion}{2.1 -- 24/11/2026}
\newcommand{\DocStatus}{Chờ duyệt}
\newcommand{\DocAuthor}{\canDien[Project Manager]}
\newcommand{\DocReviewers}{Sponsor, Product Owner, Project Manager}
\newcommand{\DocReaders}{Sponsor, Product Owner, Project Manager}
\newcommand{\CoverDate}{TP. HỒ CHÍ MINH, THÁNG 11/2026}
\hypersetup{pdftitle={SumTalk -- Hồ sơ khởi động dự án},pdfauthor={Nhóm SumTalk}}

\begin{document}

\makereportcover

% ============================================================ Front matter
\chapter*{Thông tin tài liệu}
\addcontentsline{toc}{chapter}{Thông tin tài liệu}

\docinfotable

\section*{Lịch sử phiên bản}

\begin{table}[H]
  \centering
  \begin{tabularx}{\linewidth}{@{}Q{1.6cm}Q{2.3cm}LP{3.2cm}@{}}
    \toprule
    \thead{Phiên bản} & \thead{Ngày} & \thead{Nội dung thay đổi} & \thead{Người thực hiện} \\
    \midrule
    1.0 & 12/11/2026 & Bản đầu, gồm ba tài liệu riêng: Project Charter, Business Case và lựa chọn phương pháp & \canDien \\
    2.0 & 13/11/2026 & Gộp thành một tài liệu, đánh lại số mục (CR-08) & \canDien \\
    2.1 & 24/11/2026 & Chuyển sang định dạng báo cáo PDF (LaTeX), sắp xếp lại thành năm chương & \canDien \\
    \bottomrule
  \end{tabularx}
\end{table}

\heading{Quy ước trong tài liệu}
\begin{itemize}
  \item \canDien[\ldots]: số liệu đã có ở Azure Boards, dashboard chi phí hoặc biên bản họp nhưng chưa chép
        vào tài liệu; cần bổ sung trước khi trình duyệt.
  \item \textbf{Giả định}: con số đặt ra để tính toán. Mọi bảng tài chính đều ghi rõ cách tính để thay bằng số
        thật rồi tính lại.
\end{itemize}

\tableofcontents

\chapter*{Tóm tắt}
\addcontentsline{toc}{chapter}{Tóm tắt}

Ghi biên bản họp hiện vẫn làm thủ công: thư ký vừa nghe vừa gõ và mất thêm 30--60 phút sau mỗi buổi để hoàn
thiện. Họp trực tuyến và họp trộn tiếng Anh -- tiếng Việt càng làm việc ghi chép dễ sót, trong khi các dịch vụ
ghi chú bằng AI sẵn có đưa audio lên cloud nước ngoài và tính phí theo người dùng. Dự án \textbf{SumTalk} xây
dựng một ứng dụng ghi âm, phiên âm, tóm tắt và dịch biên bản cuộc họp, chạy được trên máy của người dùng hoặc
trên server của đơn vị.

Tài liệu này là \textbf{hồ sơ khởi động dự án} (Project Initiation Document -- PID), làm căn cứ để Sponsor phê
duyệt khởi động và cấp ngân sách. Hồ sơ gồm ba phần. \emph{Điều lệ dự án} xác định sáu mục tiêu SMART (M1--M6),
phạm vi, sản phẩm bàn giao, các bên liên quan, ngân sách, mốc và phân quyền ra quyết định. \emph{Business case}
so sánh ba phương án: giữ nguyên cách làm hiện tại, tự xây SumTalk, hoặc mua dịch vụ SaaS. Với các giả định
nêu trong tài liệu, phương án tự xây đạt NPV khoảng \textbf{+109,1 triệu VND} sau 3 năm, ROI \textbf{89,4\%}
và hoàn vốn sau khoảng \textbf{16 tháng}, vượt xa phương án SaaS (NPV +25,3 triệu VND) và giữ dữ liệu họp
trong đơn vị. Cuối cùng, tài liệu lý giải việc chọn \textbf{Scrum} làm phương pháp quản lý, bổ sung các tài
liệu quản lý cần cho việc báo cáo với Sponsor.

\textbf{Từ khoá:} hồ sơ khởi động dự án, điều lệ dự án, business case, NPV, ROI, Scrum, speech-to-text,
tóm tắt cuộc họp.

\chapter*{Danh mục từ viết tắt}
\addcontentsline{toc}{chapter}{Danh mục từ viết tắt}

\begin{center}
  \setstretch{1.1}
  \begin{tabular}{@{}P{2.4cm}P{11cm}@{}}
    \toprule
    \thead{Viết tắt} & \thead{Ý nghĩa} \\
    \midrule
    API   & Application Programming Interface -- giao diện lập trình ứng dụng \\
    AWS   & Amazon Web Services (EC2: máy chủ ảo; S3: lưu trữ đối tượng) \\
    BA    & Business Analyst -- chuyên viên phân tích nghiệp vụ \\
    CPI   & Cost Performance Index -- chỉ số hiệu quả chi phí \\
    CPM   & Critical Path Method -- phương pháp đường găng \\
    DoD   & Definition of Done -- tiêu chí hoàn thành \\
    E2E   & End-to-End -- kiểm thử đầu cuối \\
    EVA   & Earned Value Analysis -- phân tích giá trị thu được \\
    LLM   & Large Language Model -- mô hình ngôn ngữ lớn \\
    NPV   & Net Present Value -- giá trị hiện tại ròng \\
    PID   & Project Initiation Document -- hồ sơ khởi động dự án \\
    PM    & Project Manager -- quản lý dự án \\
    PO    & Product Owner -- chủ sản phẩm \\
    POC   & Proof of Concept -- bản thử nghiệm khái niệm \\
    QC    & Quality Control -- kiểm soát chất lượng \\
    ROI   & Return on Investment -- tỷ suất hoàn vốn \\
    SaaS  & Software as a Service -- phần mềm dạng dịch vụ \\
    SLA   & Service Level Agreement -- cam kết mức dịch vụ \\
    SMART & Specific, Measurable, Achievable, Relevant, Time-bound \\
    STT   & Speech-to-Text -- chuyển giọng nói thành văn bản \\
    UAT   & User Acceptance Testing -- kiểm thử chấp nhận của người dùng \\
    \bottomrule
  \end{tabular}
\end{center}

\cleardoublepage\phantomsection
\addcontentsline{toc}{chapter}{\listfigurename}
\listoffigures

\cleardoublepage\phantomsection
\addcontentsline{toc}{chapter}{\listtablename}
\listoftables

% ============================================================ Chapter 1
\chapter{Giới thiệu}
\label{chap:intro}

\section{Bối cảnh và vấn đề}

Trong nhiều đơn vị, biên bản họp vẫn được ghi bằng tay. Thư ký vừa theo dõi nội dung vừa gõ, rồi mất thêm
30--60 phút sau mỗi buổi để chỉnh sửa và gửi đi. Ba xu hướng đang làm việc này khó hơn:
\begin{itemize}
  \item \textbf{Họp trực tuyến} trên Meet, Teams hay Zoom: âm thanh nằm trong cuộc gọi, muốn nghe lại phải ghi
        âm riêng.
  \item \textbf{Họp trộn ngôn ngữ}: người nói chuyển qua lại giữa tiếng Anh và tiếng Việt, nên rất dễ ghi sót
        hoặc phải dịch tay.
  \item \textbf{Dịch vụ AI ghi chú có sẵn} thường đưa audio lên cloud nước ngoài, hỗ trợ tiếng Việt và họp đa
        ngôn ngữ còn yếu, lại tính phí theo từng người dùng.
\end{itemize}
Nói cách khác, vấn đề không chỉ là tốn thời gian ghi chép, mà còn là \emph{chất lượng} biên bản (sót quyết định,
sót việc cần làm) và \emph{chủ quyền dữ liệu} của các cuộc họp có nội dung nhạy cảm.

\section{Giải pháp đề xuất: SumTalk}

SumTalk là ứng dụng ghi âm, phiên âm, tóm tắt và dịch biên bản cuộc họp. Người dùng ghi âm trực tiếp (micro,
hoặc micro cùng âm thanh cuộc họp trực tuyến) hay tải lên file có sẵn; hệ thống dùng Whisper để chuyển giọng
nói thành văn bản, giữ đúng ngôn ngữ của từng câu, hiển thị transcript ngay khi đang nói, sau đó dùng LLM để
tạo biên bản tóm tắt và dịch Việt $\leftrightarrow$ Anh. Sản phẩm có bản web triển khai trên AWS và ứng dụng
desktop cho Windows, nên có thể chạy hoàn toàn trên máy hoặc server của đơn vị. Phạm vi chi tiết được nêu ở
Mục~\ref{sec:scope}.

\section{Mục đích và người đọc của tài liệu}

Tài liệu trả lời năm câu hỏi trước khi dự án được khởi động: \emph{vì sao} làm SumTalk, \emph{nhắm tới} mục
tiêu gì, làm trong \emph{phạm vi} nào, có \emph{đáng đầu tư} không, và \emph{quản lý} theo phương pháp nào.
Người đọc chính là Sponsor (phê duyệt và cấp ngân sách), Product Owner và Project Manager. Các kế hoạch chi
tiết (WBS, lịch trình, rủi ro, chi phí) nằm trong \docref{03}{Kế hoạch quản lý dự án}~\cite{pmp}; yêu cầu
sản phẩm nằm trong \docref{02}{Đặc tả yêu cầu phần mềm}~\cite{srs}.

\section{Cấu trúc tài liệu}

Phần còn lại của tài liệu gồm bốn chương:
\begin{itemize}
  \item \textbf{Chương~\ref{chap:charter} -- Điều lệ dự án:} mục tiêu, phạm vi, sản phẩm bàn giao, các bên liên
        quan, ngân sách, mốc chính, giả định, rủi ro cấp cao và phân quyền.
  \item \textbf{Chương~\ref{chap:business} -- Business case:} bài toán nghiệp vụ, các phương án, giả định tài
        chính, dòng tiền, NPV, ROI, thời gian hoàn vốn và phân tích độ nhạy.
  \item \textbf{Chương~\ref{chap:method} -- Phương pháp quản lý dự án:} so sánh Waterfall, Scrum, Kanban và mô
        hình lai; cách áp dụng Scrum trong dự án.
  \item \textbf{Chương~\ref{chap:conclusion} -- Kết luận và đề nghị phê duyệt.}
\end{itemize}

% ============================================================ Chapter 2
\chapter{Điều lệ dự án}
\label{chap:charter}

Điều lệ dự án (Project Charter) chính thức xác lập dự án SumTalk, trao quyền cho Project Manager sử dụng nguồn
lực đã duyệt và là điểm tham chiếu khi có yêu cầu thay đổi phạm vi. Thông tin chính được tóm tắt trong
Bảng~\ref{tab:charter-info}.

\begin{table}[htbp]
  \centering
  \caption{Thông tin chính của dự án}
  \label{tab:charter-info}
  \begin{tabularx}{0.9\linewidth}{@{}P{5cm}L@{}}
    \toprule
    Tên dự án & SumTalk -- ghi âm, phiên âm, tóm tắt và dịch biên bản cuộc họp \\
    Nhà tài trợ (Sponsor) & \canDien[đại diện đơn vị tài trợ] \\
    Product Owner & \canDien \\
    Project Manager / Scrum Master & \canDien \\
    Ngày phê duyệt & \canDien \\
    Phiên bản điều lệ & 1.0 -- 12/11/2026 \\
    \bottomrule
  \end{tabularx}
\end{table}

\section{Mục tiêu dự án}

Mỗi mục tiêu được viết theo tiêu chí SMART: có thước đo kiểm chứng được và gắn với một sprint hoặc mốc cụ thể.
Các mục tiêu M1--M4 thuộc về sản phẩm, M5 về khả năng triển khai, còn M6 ràng buộc ngân sách
(Bảng~\ref{tab:objectives}). Một số thước đo đã có kết quả đo thực tế, được ghi kèm để đối chiếu.

\begin{table}[htbp]
  \centering
  \caption{Mục tiêu dự án theo tiêu chí SMART}
  \label{tab:objectives}
  \small
  \begin{tabularx}{\linewidth}{@{}P{0.9cm}LLP{2.1cm}@{}}
    \toprule
    \thead{Mã} & \thead{Mục tiêu} & \thead{Thước đo} & \thead{Hạn} \\
    \midrule
    M1 & Ghi âm hoặc tải file cuộc họp và nhận transcript đúng ngôn ngữ người nói (Anh, Việt, trộn)
       & Chạy được trên Chrome/Edge, điện thoại, app Windows; test E2E đạt & Sprint 1--2 \\
    M2 & Transcript hiện trực tiếp khi đang nói
       & Chữ hiện trễ $\leq$ 3--5\,s (đo được: khoảng 1\,s cho từ mới, 1,9\,s cho câu đã chốt) & Sprint 4 \\
    M3 & Tự động tóm tắt (tổng quan, ý chính, quyết định, việc cần làm) và dịch Việt $\leftrightarrow$ Anh
       & Story \#67 và \#60 đạt toàn bộ acceptance criteria & Sprint 3 \\
    M4 & Xử lý được cuộc họp từ 60 phút trở lên
       & Đã đo: file 11,6 phút xong trong 18--57\,s, dòng đầu hiện sau 3--5\,s & Sprint 4 \\
    M5 & Có bản web trên AWS để test, demo và bộ cài Windows không bắt tải model nặng
       & Có URL HTTPS và bộ cài \texttt{SumTalk-Setup-x.y.z.exe} & Sprint 4 \\
    M6 & Nằm trong ngân sách
       & Chi phí thực tế $\leq$ ngân sách duyệt (Mục~\ref{sec:budget}), CPI $\geq$ 0,9 & Cuối dự án \\
    \bottomrule
  \end{tabularx}
\end{table}

\section{Phạm vi}
\label{sec:scope}

\heading{Trong phạm vi}
\begin{itemize}
  \item \textbf{Ghi âm:} micro, hoặc micro cùng âm thanh cuộc họp trực tuyến; tạm dừng và tiếp tục; vẫn ghi khi
        mất mạng rồi tự đồng bộ khi có mạng lại.
  \item \textbf{Tải file:} file audio/video đến 1\,GB, nhiều định dạng, tự chuẩn hoá về WAV 16\,kHz.
  \item \textbf{Speech-to-text:} Whisper chạy local (GPU/CPU) hoặc qua API/server; giữ ngôn ngữ gốc của từng câu.
  \item \textbf{Transcript:} hiển thị trực tiếp trong khi họp và bản cuối chính xác; tìm kiếm, phát lại đồng bộ,
        xuất TXT, SRT, VTT, JSON và WAV.
  \item \textbf{Tóm tắt và dịch} bằng LLM (tiếng Việt, tiếng Anh), xuất được bản dịch.
  \item \textbf{Nền tảng:} ứng dụng desktop Windows (Electron), bản web trên AWS (EC2 và S3), cùng bộ tài liệu dự án.
\end{itemize}

\heading{Ngoài phạm vi (phiên bản này)}
\begin{itemize}
  \item Tài khoản người dùng, phân quyền, đa tổ chức; bản test dùng một mật khẩu chung.
  \item Nhận diện người nói (diarization), dịch trực tiếp khi đang họp (\#84), trang việc cần làm (\#79).
  \item Ứng dụng iOS/Android native, ghi âm khi khoá màn hình trên iOS.
  \item Cam kết SLA vận hành thương mại.
\end{itemize}

\section{Sản phẩm bàn giao}

Bảng~\ref{tab:deliverables} liệt kê năm nhóm sản phẩm bàn giao. Mã nguồn được tổ chức theo nhánh feature của
từng sprint; bản web trên AWS và ứng dụng desktop thuộc mốc M4 (Mục~\ref{sec:milestones}).

\begin{table}[htbp]
  \centering
  \caption{Sản phẩm bàn giao chính}
  \label{tab:deliverables}
  \begin{tabularx}{\linewidth}{@{}P{3.6cm}L@{}}
    \toprule
    \thead{Sản phẩm} & \thead{Mô tả} \\
    \midrule
    Mã nguồn & Kho \texttt{DatTapTanhCode/SumTalk}, các nhánh feature theo sprint \\
    Bản web & Docker image và hướng dẫn deploy trên AWS EC2/S3 \\
    Ứng dụng desktop & Bộ cài Windows (NSIS) và bản GPU dạng thư mục \\
    Tài liệu & README, ARCHITECTURE, bộ sáu tài liệu quản lý dự án, test case \\
    Kết quả kiểm thử & pytest (backend), Playwright (trình duyệt, desktop) và báo cáo kiểm thử~\cite{qtd} \\
    \bottomrule
  \end{tabularx}
\end{table}

\section{Các bên liên quan}

Ba bên có ảnh hưởng cao là Sponsor, Product Owner và nhóm phát triển; họ tham gia trực tiếp vào các sự kiện của
mỗi sprint. Các bên còn lại được cập nhật qua kênh phù hợp với mức ảnh hưởng (Bảng~\ref{tab:stakeholders}).
Phân công chi tiết theo ma trận RACI được trình bày trong~\cite{pmp}.

\begin{table}[htbp]
  \centering
  \caption{Các bên liên quan và cách làm việc}
  \label{tab:stakeholders}
  \small
  \begin{tabularx}{\linewidth}{@{}P{3.4cm}LQ{1.9cm}L@{}}
    \toprule
    \thead{Bên liên quan} & \thead{Vai trò / mối quan tâm} & \thead{Ảnh hưởng} & \thead{Cách làm việc} \\
    \midrule
    Sponsor (đơn vị tài trợ) & Cấp ngân sách, duyệt phạm vi và mốc, nghiệm thu & Cao
      & Sprint Review, báo cáo trạng thái mỗi sprint \\
    Product Owner & Ưu tiên backlog, chấp nhận story & Cao & Sprint Planning/Review, Azure Boards \\
    Nhóm phát triển & Xây dựng sản phẩm, kiểm thử đơn vị & Cao & Daily, code review \\
    BA & Yêu cầu, acceptance criteria, mẫu tham chiếu để đánh giá & Trung bình & Refinement \\
    QC & Test case, test E2E, báo lỗi & Trung bình & Test plan, Bug trên Azure Boards \\
    Người dùng thử (thư ký, trưởng nhóm họp) & Dễ dùng, đúng nội dung, bảo mật & Trung bình & UAT trên bản web \\
    Nhà cung cấp (AWS, Groq/OpenAI, Ollama) & Hạ tầng, API & Thấp & Theo dõi chi phí, rate limit \\
    \bottomrule
  \end{tabularx}
\end{table}

\section{Ngân sách tóm tắt}
\label{sec:budget}

Chi phí chủ yếu là nhân công, tính theo đơn giá 100.000~VND/giờ:
\begin{equation}
  \label{eq:labour}
  \text{Chi phí nhân công} = \text{tổng giờ kế hoạch} \times 100.000~\text{VND}
\end{equation}
Chi phí ngoài nhân công (AWS, API) ước tính khoảng 3~triệu VND/tháng khi chạy bản test 24/7. Ngân sách giữ thêm
10\% dự phòng (Bảng~\ref{tab:budget}). Kế hoạch chi phí chi tiết theo sprint nằm trong~\cite{pmp}.

\begin{table}[htbp]
  \centering
  \caption{Ngân sách tóm tắt}
  \label{tab:budget}
  \begin{tabularx}{\linewidth}{@{}P{5.2cm}LP{3.3cm}@{}}
    \toprule
    \thead{Khoản} & \thead{Giá trị} & \thead{Nguồn} \\
    \midrule
    Đơn giá nhân công & 100.000 VND/giờ & Dashboard chi phí \\
    Tổng giờ kế hoạch & \canDien[tổng giờ theo sprint] & Azure Boards \\
    Chi phí nhân công kế hoạch & Theo công thức~\eqref{eq:labour} & -- \\
    Chi phí ngoài nhân công & Khoảng 3 triệu VND/tháng khi chạy bản test 24/7 & Ước tính \\
    Dự phòng & 10\% tổng ngân sách & -- \\
    \bottomrule
  \end{tabularx}
\end{table}

\section{Mốc chính}
\label{sec:milestones}

Dự án có sáu mốc: M0 khi được duyệt khởi động, M1--M4 ứng với kết quả cuối Sprint 1--4, và M5 khi nghiệm thu
(Hình~\ref{fig:milestones}). Ngày cụ thể của từng mốc
được ghi trong Bảng~\ref{tab:milestones}; lịch chi tiết với sơ đồ mạng, CPM và Gantt nằm trong~\cite{pmp}.

\begin{figure}[htbp]
  \centering
  \begin{tikzpicture}[
      ms/.style={circle,draw=white,line width=1pt,fill=primaryblue,inner sep=3pt},
      lbl/.style={text width=3cm,align=center,font=\footnotesize,inner sep=1pt,
                  execute at begin node=\setstretch{1.05}}]
    \draw[-{Stealth[length=3mm]},line width=1.2pt,primaryblue] (0,0) -- (15.8,0);
    \foreach \x/\m/\when/\what in {
        1.5/M0/Khởi động/{Charter, backlog, yêu cầu và UI/UX được duyệt},
        6.6/M2/Hết Sprint 2/{Benchmark và chọn model LLM},
        11.7/M4/Hết Sprint 4/{Desktop, transcript streaming, STT qua server, bản web AWS}} {
      \node[ms] at (\x,0) {};
      \node[lbl,above=7pt] at (\x,0) {\textbf{\m} -- \when\\ \what};
    }
    \foreach \x/\m/\when/\what in {
        4.05/M1/Hết Sprint 1/{Ghi âm, upload, transcript gốc (POC E1)},
        9.15/M3/Hết Sprint 3/{Tóm tắt (\#67), dịch (\#60), test E2E},
        14.25/M5/Kết thúc/{Nghiệm thu, báo cáo tổng kết}} {
      \node[ms] at (\x,0) {};
      \node[lbl,below=7pt] at (\x,0) {\textbf{\m} -- \when\\ \what};
    }
  \end{tikzpicture}
  \caption{Các mốc chính của dự án}
  \label{fig:milestones}
\end{figure}

\begin{table}[htbp]
  \centering
  \caption{Mốc chính và ngày hoàn thành}
  \label{tab:milestones}
  \begin{tabularx}{\linewidth}{@{}P{3.2cm}LP{2.6cm}@{}}
    \toprule
    \thead{Mốc} & \thead{Nội dung} & \thead{Ngày} \\
    \midrule
    M0 & Charter, backlog, tài liệu yêu cầu và UI/UX được duyệt & \canDien \\
    M1 (hết Sprint 1) & Ghi âm, upload, transcript gốc chạy được (POC E1) & \canDien \\
    M2 (hết Sprint 2) & Benchmark và chọn model LLM & \canDien \\
    M3 (hết Sprint 3) & Tóm tắt (\#67), dịch (\#60), test E2E (\#99, \#100) & \canDien \\
    M4 (hết Sprint 4) & App desktop, transcript streaming, STT qua server, bản web trên AWS & \canDien \\
    M5 & Nghiệm thu, báo cáo tổng kết & \canDien \\
    \bottomrule
  \end{tabularx}
\end{table}

\section{Giả định và ràng buộc}

\begin{itemize}
  \item \textbf{Nhân lực:} nhóm làm bán thời gian, song song với công việc khác; mỗi sprint dài
        \canDien[1--2] tuần.
  \item \textbf{Hạ tầng:} không có ngân sách thuê GPU cố định. Nhóm phát triển dùng laptop RTX 4050; bản test
        trên AWS dùng EC2 CPU kết hợp Whisper API, hoặc bật GPU theo giờ.
  \item \textbf{Bảo mật dữ liệu:} nội dung họp có thể nhạy cảm, nên mặc định chạy local; khi deploy bắt buộc có
        HTTPS, đăng nhập và bucket S3 private.
  \item \textbf{Công cụ:} Azure Boards (backlog), GitHub (mã nguồn), Python/FastAPI, React, Electron.
\end{itemize}

\section{Rủi ro cấp cao}

Bốn rủi ro sau được xác định ngay khi khởi động; phân tích xác suất, tác động và kế hoạch ứng phó nằm trong
Risk Register~\cite{pmp}.
\begin{enumerate}
  \item Độ chính xác của STT với tiếng Việt và họp trộn ngôn ngữ.
  \item Phụ thuộc GPU và API bên ngoài: chi phí và giới hạn số lượng yêu cầu (rate limit).
  \item Trễ tiến độ do nhóm làm bán thời gian.
  \item Bảo mật dữ liệu cuộc họp.
\end{enumerate}

\section{Phân quyền ra quyết định}

Bảng~\ref{tab:decision-rights} xác định ai có quyền quyết định cuối cùng (in đậm) với từng loại quyết định. Quy
tắc quan trọng nhất là thay đổi phạm vi, mốc hoặc ngân sách vượt 10\% phải được Sponsor phê duyệt; mọi thay đổi
được ghi vào Change Log~\cite{mcr}.

% Last table of the chapter: [H] so it is not pushed onto a page of its own.
\begin{table}[H]
  \centering
  \caption{Phân quyền ra quyết định}
  \label{tab:decision-rights}
  \small
  \begin{tabularx}{\linewidth}{@{}LQ{1.9cm}Q{1.9cm}Q{1.9cm}Q{1.9cm}@{}}
    \toprule
    \thead{Quyết định} & \thead{PM} & \thead{PO} & \thead{Nhóm Dev} & \thead{Sponsor} \\
    \midrule
    Thêm/bớt story trong sprint đang chạy & đề xuất & \textbf{quyết định} & tư vấn & -- \\
    Thay đổi phạm vi, mốc, ngân sách $>$ 10\% & đề xuất & đồng ý & -- & \textbf{phê duyệt} \\
    Giải pháp kỹ thuật & tư vấn & -- & \textbf{quyết định} & -- \\
    Phân công, lịch họp & \textbf{quyết định} & -- & tư vấn & -- \\
    Chấp nhận story (Definition of Done) & -- & \textbf{quyết định} & -- & -- \\
    \bottomrule
  \end{tabularx}
\end{table}

% ============================================================ Chapter 3
\chapter{Business case}
\label{chap:business}

Chương này đánh giá dự án có đáng đầu tư hay không. \textbf{Kết luận: nên tự xây SumTalk.} Với các giả định ở
Mục~\ref{sec:fin-assumptions}, sau 3 năm phương án tự xây đạt NPV khoảng +109,1 triệu VND, ROI 89,4\% và hoàn
vốn sau khoảng 1,3 năm. Phương án mua SaaS chỉ đạt NPV khoảng +25,3 triệu VND, ROI 12,8\%, và audio họp phải đưa
ra cloud nước ngoài.

Mọi số tiền trong chương là \textbf{ước tính dựa trên giả định} với tỷ giá 26.000~VND/USD. Khi có số thật (tổng
giờ công trên dashboard, số cuộc họp của đơn vị dùng thử), có thể tính lại bằng công cụ
\texttt{docs/project/tools/pm\_calc.py}~\cite{pmcalc}.

\section{Bài toán nghiệp vụ}

Bảng~\ref{tab:as-is-to-be} so sánh hiện trạng với trạng thái mong muốn ở bốn khía cạnh. Ba khía cạnh đầu liên
quan tới năng suất; khía cạnh thứ tư là lý do khiến các dịch vụ AI có sẵn khó được chấp nhận.

\begin{table}[htbp]
  \centering
  \caption{Hiện trạng và trạng thái mong muốn}
  \label{tab:as-is-to-be}
  \begin{tabularx}{\linewidth}{@{}P{2.8cm}LL@{}}
    \toprule
    \thead{Khía cạnh} & \thead{Hiện trạng} & \thead{Mong muốn} \\
    \midrule
    Ghi biên bản & Thư ký gõ tay trong và sau họp, mất 30--60 phút mỗi buổi
      & Transcript có sẵn khi họp xong; biên bản (tóm tắt, quyết định, việc cần làm) tự tạo \\
    Họp trực tuyến & Phải ghi âm riêng rồi nghe lại & Thu cả micro và âm thanh cuộc gọi, chữ hiện ngay khi nói \\
    Họp trộn Anh -- Việt & Dễ sót, phải dịch tay & Giữ đúng ngôn ngữ từng câu, dịch Việt $\leftrightarrow$ Anh một chạm \\
    Bảo mật & Dịch vụ AI bên ngoài đưa dữ liệu ra nước ngoài & Chạy trên máy hoặc server của đơn vị; bucket S3 private \\
    \bottomrule
  \end{tabularx}
\end{table}

\section{Giá trị mang lại}

\begin{itemize}
  \item \textbf{Tiết kiệm thời gian:} giả định tiết kiệm 45 phút cho mỗi cuộc họp cần biên bản; đây là nguồn
        lợi ích duy nhất được lượng hoá trong Mục~\ref{sec:cashflow}.
  \item \textbf{Chất lượng biên bản:} không sót quyết định hay việc cần làm; tìm kiếm và phát lại đúng đoạn.
  \item \textbf{Khả năng tiếp cận:} người vắng mặt đọc tóm tắt hoặc bản dịch thay vì nghe lại một giờ ghi âm.
  \item \textbf{Chủ quyền dữ liệu:} tự host được, trên máy cá nhân hoặc server riêng.
\end{itemize}

\section{Các phương án}

Ba phương án được đưa vào so sánh (Bảng~\ref{tab:options}). Phương án P0 là mốc tham chiếu; P1 và P2 mang lại
cùng lợi ích về thời gian nhưng khác nhau về chi phí và nơi lưu dữ liệu.

\begin{table}[htbp]
  \centering
  \caption{Các phương án được xem xét}
  \label{tab:options}
  \small
  \begin{tabularx}{\linewidth}{@{}P{2.7cm}LLL@{}}
    \toprule
    \thead{Phương án} & \thead{Mô tả} & \thead{Ưu điểm} & \thead{Nhược điểm} \\
    \midrule
    P0 -- Giữ nguyên & Ghi tay như hiện tại & Không tốn đầu tư & Tốn giờ công, sót nội dung \\
    \textbf{P1 -- Tự xây SumTalk} & Dự án này: web và desktop, Whisper và LLM
      & Làm theo nhu cầu tiếng Việt và họp trộn ngôn ngữ, tự host, chi phí vận hành thấp
      & Tốn công phát triển, cần bảo trì \\
    P2 -- Mua SaaS & Dịch vụ AI note-taker trả phí theo người & Dùng được ngay
      & Phí theo người cao, dữ liệu ra nước ngoài, tiếng Việt và họp trộn còn yếu \\
    \bottomrule
  \end{tabularx}
\end{table}

\section{Giả định tài chính}
\label{sec:fin-assumptions}

Bảng~\ref{tab:fin-assumptions} liệt kê các giả định đầu vào. Lợi ích tối đa mỗi năm bằng số cuộc họp nhân với
thời gian tiết kiệm và đơn giá giờ công: $1.920 \times 0{,}75 \times 100.000 = 144$ triệu VND, nhân với tỷ lệ
áp dụng của từng năm.

\begin{table}[htbp]
  \centering
  \caption{Giả định tài chính}
  \label{tab:fin-assumptions}
  \small
  \begin{tabularx}{\linewidth}{@{}P{3.6cm}LP{4.3cm}@{}}
    \toprule
    \thead{Giả định} & \thead{Giá trị} & \thead{Ghi chú} \\
    \midrule
    Tổng công phát triển & \textbf{600 giờ} $\times$ 100.000 VND = \textbf{60.000.000 VND}
      & \canDien[thay bằng tổng giờ thực tế] \\
    Đơn vị dùng thử & 40 cuộc họp/tuần $\times$ 48 tuần = 1.920 cuộc/năm & Ví dụ một công ty khoảng 50 người \\
    Thời gian tiết kiệm & 0,75 giờ/cuộc họp, quy đổi 100.000 VND/giờ & Lợi ích tối đa 144.000.000 VND/năm \\
    Tỷ lệ áp dụng & Năm 1: 50\%, năm 2: 80\%, năm 3: 90\% & -- \\
    Vận hành P1 & 35.744.000 VND/năm & Chi tiết ở Bảng~\ref{tab:opex} \\
    Phí SaaS P2 & 20 người $\times$ 15 USD/tháng = 93.600.000 VND/năm & Giá tham khảo, cần kiểm tra \\
    Suất chiết khấu, thời gian & 10\%/năm, 3 năm & -- \\
    \bottomrule
  \end{tabularx}
\end{table}

Chi phí vận hành của P1 chủ yếu là máy chủ EC2; chi phí API cho STT và LLM chỉ chiếm khoảng 9\% tổng
(Bảng~\ref{tab:opex}).

\begin{table}[htbp]
  \centering
  \caption{Chi phí vận hành hằng năm của phương án P1}
  \label{tab:opex}
  \begin{tabular}{@{}P{10cm}r@{}}
    \toprule
    \thead{Khoản} & \thead{VND/năm} \\
    \midrule
    EC2 t3.large, EBS và IP (khoảng 85 USD/tháng) & 26.520.000 \\
    Whisper API (1.600 giờ audio $\times$ 0,04 USD) & 1.664.000 \\
    LLM API tóm tắt và dịch (khoảng 5 USD/tháng) & 1.560.000 \\
    Bảo trì (10\% chi phí phát triển) & 6.000.000 \\
    \midrule
    \textbf{Tổng} & \textbf{35.744.000} \\
    \bottomrule
  \end{tabular}
\end{table}

\section{Dòng tiền, NPV, ROI và thời gian hoàn vốn}
\label{sec:cashflow}

Dòng tiền ròng mỗi năm bằng lợi ích trừ chi phí. Bảng~\ref{tab:cashflow} trình bày dòng tiền của hai phương án
có đầu tư; lợi ích của P1 và P2 như nhau vì cùng tiết kiệm thời gian ghi biên bản.

\begin{table}[htbp]
  \centering
  \caption{Dòng tiền theo năm (triệu VND)}
  \label{tab:cashflow}
  \begin{tabular}{@{}lrrrr@{}}
    \toprule
    \thead{Hạng mục} & \thead{Năm 0} & \thead{Năm 1} & \thead{Năm 2} & \thead{Năm 3} \\
    \midrule
    Lợi ích (P1 và P2) & 0 & 72,0 & 115,2 & 129,6 \\
    \midrule
    P1 -- chi phí & 60,0 & 35,7 & 35,7 & 35,7 \\
    P1 -- dòng tiền ròng & $-60{,}0$ & 36,3 & 79,5 & 93,9 \\
    P1 -- luỹ kế & $-60{,}0$ & $-23{,}7$ & 55,7 & 149,6 \\
    \midrule
    P2 -- chi phí & 0 & 93,6 & 93,6 & 93,6 \\
    P2 -- dòng tiền ròng & 0 & $-21{,}6$ & 21,6 & 36,0 \\
    \bottomrule
  \end{tabular}
\end{table}

Ba chỉ số được dùng để so sánh. Giá trị hiện tại ròng chiết khấu dòng tiền $CF_t$ của từng năm với suất
chiết khấu $r$:
\begin{equation}
  \label{eq:npv}
  \mathrm{NPV} = \sum_{t=0}^{3} \frac{CF_t}{(1+r)^t}, \qquad r = 10\%
\end{equation}
Áp dụng cho P1 với dòng tiền chưa làm tròn:
\begin{equation*}
  \mathrm{NPV}_{P1} = -60 + \frac{36{,}256}{1{,}1} + \frac{79{,}456}{1{,}1^2} + \frac{93{,}856}{1{,}1^3}
  \approx 109{,}1~\text{triệu VND}
\end{equation*}
Tỷ suất hoàn vốn so sánh tổng lợi ích $B_t$ với tổng chi phí $C_t$ trong 3 năm:
\begin{equation}
  \label{eq:roi}
  \mathrm{ROI} = \frac{\sum_t B_t - \sum_t C_t}{\sum_t C_t}
  \qquad\Rightarrow\qquad
  \mathrm{ROI}_{P1} = \frac{316{,}8 - 167{,}2}{167{,}2} \approx 89{,}4\%
\end{equation}
Thời gian hoàn vốn được nội suy trong năm dòng tiền luỹ kế đổi dấu. Với P1, luỹ kế hết năm 1 là
$-23{,}7$ triệu và năm 2 thu thêm 79,5 triệu:
\begin{equation}
  \label{eq:payback}
  T_{\text{hoàn vốn}} = 1 + \frac{23{,}7}{79{,}5} \approx 1{,}30~\text{năm} \quad (\text{khoảng 16 tháng})
\end{equation}

Kết quả của cả ba phương án được tổng hợp trong Bảng~\ref{tab:fin-results}, và dòng tiền luỹ kế được minh hoạ
ở Hình~\ref{fig:cumulative}. P1 lỗ nhiều hơn trong năm đầu vì phải trả chi phí phát triển, nhưng từ năm 2 trở đi
vượt hẳn P2 do chi phí vận hành thấp hơn phí SaaS gần 58 triệu VND mỗi năm.

\begin{table}[htbp]
  \centering
  \caption{So sánh chỉ số tài chính của các phương án}
  \label{tab:fin-results}
  \begin{tabularx}{\linewidth}{@{}LCCC@{}}
    \toprule
    \thead{Chỉ số} & \thead{P0 -- Giữ nguyên} & \thead{P1 -- Tự xây} & \thead{P2 -- SaaS} \\
    \midrule
    NPV (3 năm, 10\%) & 0 & \textbf{+109,1 triệu} & +25,3 triệu \\
    ROI (3 năm) & -- & \textbf{89,4\%} & 12,8\% \\
    Thời gian hoàn vốn & -- & \textbf{1,3 năm (khoảng 16 tháng)} & khoảng 2 năm \\
    Nơi lưu dữ liệu họp & Trong đơn vị & Trong đơn vị & Nước ngoài \\
    \bottomrule
  \end{tabularx}
\end{table}

\begin{figure}[htbp]
  \centering
  \begin{tikzpicture}
    \begin{axis}[
        width=0.88\linewidth, height=7.2cm,
        xmin=-0.25, xmax=3.35, ymin=-85, ymax=175,
        xtick={0,1,2,3}, xticklabels={Năm 0,Năm 1,Năm 2,Năm 3},
        ylabel={Dòng tiền luỹ kế (triệu VND)},
        grid=major, grid style={gray!25},
        legend pos=north west, legend cell align=left,
        legend style={font=\small,draw=gray!40},
        tick label style={font=\small},
        /pgf/number format/.cd, use comma, 1000 sep={.}]
      \addplot[gray!70, forget plot] coordinates {(-0.25,0) (3.35,0)};
      \addplot[very thick, primaryblue, mark=*,
               nodes near coords, every node near coord/.append style={font=\scriptsize,anchor=north west}]
        coordinates {(0,-60) (1,-23.7) (2,55.7) (3,149.6)};
      \addlegendentry{P1 -- Tự xây SumTalk}
      \addplot[thick, chartgray, dashed, mark=square*]
        coordinates {(0,0) (1,-21.6) (2,0) (3,36)};
      \addlegendentry{P2 -- Mua SaaS}
      \addplot[only marks, mark=diamond*, mark size=3.5pt, todored]
        coordinates {(1.30,0)};
      \addlegendentry{Điểm hoàn vốn của P1 ($\approx$ 1,3 năm)}
    \end{axis}
  \end{tikzpicture}
  \caption{Dòng tiền luỹ kế của hai phương án có đầu tư}
  \label{fig:cumulative}
\end{figure}

\section{Phân tích độ nhạy}

Để kiểm tra mức độ chắc chắn của kết luận, NPV của P1 được tính lại khi thay đổi từng giả định chính
(Bảng~\ref{tab:sensitivity}).

\begin{table}[htbp]
  \centering
  \caption{Độ nhạy của NPV phương án P1}
  \label{tab:sensitivity}
  \begin{tabularx}{\linewidth}{@{}LP{5.4cm}@{}}
    \toprule
    \thead{Kịch bản} & \thead{NPV của P1} \\
    \midrule
    Cơ sở & +109,1 triệu \\
    Công phát triển gấp đôi (1.200 giờ) & +49,1 triệu (vẫn có lời) \\
    Lợi ích chỉ bằng một nửa (ít họp hoặc áp dụng chậm) & $-19{,}9$ triệu (lỗ) \\
    Điểm hoà vốn & Khoảng \textbf{23 cuộc họp/tuần} thì NPV = 0 \\
    \bottomrule
  \end{tabularx}
\end{table}

Kết quả nhạy với \emph{số cuộc họp thực tế} hơn là với chi phí phát triển: tăng gấp đôi công phát triển vẫn có
lời, nhưng lợi ích giảm một nửa thì lỗ. Vì vậy, trước khi mở rộng cần đo số cuộc họp và thời gian tiết kiệm thật
trong đợt UAT.

\section{Rủi ro tài chính}

\begin{itemize}
  \item \textbf{Chi phí API vượt dự kiến} khi dùng nhiều (rủi ro R6). Biện pháp: đã có cấu hình
        \texttt{LIVE\_INTERVAL\_SEC} và nhật ký token, chi phí cho từng lần gọi.
  \item \textbf{GPU đắt nếu chạy 24/7.} Biện pháp: bật/tắt theo giờ, hoặc dùng cấu hình CPU kết hợp API.
  \item \textbf{Tỷ lệ áp dụng thấp}, kéo lợi ích xuống dưới điểm hoà vốn (Bảng~\ref{tab:sensitivity}).
\end{itemize}
Các rủi ro này được theo dõi cùng các rủi ro khác trong Risk Register~\cite{pmp}.

\section{Khuyến nghị}

Chọn \textbf{phương án P1}. Bắt đầu bằng bản web trên EC2 CPU kết hợp Whisper API để giữ chi phí thấp. Chỉ chuyển
sang server GPU (\texttt{STT\_SERVER\_API\_KEYS}) khi số cuộc họp vượt khoảng 25 cuộc mỗi tuần, hoặc khi có
nhiều người dùng ứng dụng desktop.

% ============================================================ Chapter 4
\chapter{Phương pháp quản lý dự án}
\label{chap:method}

\textbf{Quyết định: dùng Scrum}, với sprint dài \canDien[1--2] tuần và backlog trên Azure Boards. Lý do chính là
các yêu cầu về AI (độ chính xác STT, chất lượng tóm tắt) chỉ biết được sau khi thử, và phạm vi đã thay đổi nhiều
lần: dời story \#60 vào Sprint 3, thêm app desktop, transcript streaming và deploy cloud. Waterfall sẽ khoá yêu
cầu quá sớm.

\section{Đặc điểm dự án ảnh hưởng tới lựa chọn}

Bảng~\ref{tab:characteristics} đối chiếu sáu đặc điểm của dự án với cách tiếp cận phù hợp cho từng đặc điểm.

\begin{table}[htbp]
  \centering
  \caption{Đặc điểm của dự án và phương pháp phù hợp}
  \label{tab:characteristics}
  \small
  \begin{tabularx}{\linewidth}{@{}P{3.1cm}LL@{}}
    \toprule
    \thead{Đặc điểm} & \thead{Ở dự án SumTalk} & \thead{Cách tiếp cận phù hợp} \\
    \midrule
    Độ rõ của yêu cầu & Rõ ở mức tính năng, mơ hồ ở mức chất lượng AI (độ chính xác, độ trễ) & Lặp, đo, điều chỉnh \\
    Rủi ro công nghệ & Cao: Whisper, LLM, GPU, trình duyệt, Electron & Làm POC sớm, demo mỗi sprint \\
    Thay đổi phạm vi & Thường xuyên (xem Change Log~\cite{mcr}) & Ưu tiên lại backlog mỗi sprint \\
    Nhóm & Nhỏ, bán thời gian, kiêm nhiều vai & Tự tổ chức, họp ngắn \\
    Sponsor & Muốn thấy tiến độ đều và sản phẩm chạy được & Sprint Review có demo \\
    Ràng buộc báo cáo & Sponsor cần báo cáo ngân sách, tiến độ, rủi ro
      & Scrum kèm bộ tài liệu quản lý tối thiểu, cập nhật theo sprint \\
    \bottomrule
  \end{tabularx}
\end{table}

Phần lớn đặc điểm của dự án (yêu cầu chưa rõ ở mức chất lượng, rủi ro công nghệ cao, phạm vi hay thay đổi) đều
nghiêng về cách làm lặp. Đặc điểm duy nhất kéo về phía kế hoạch cố định
là yêu cầu báo cáo của Sponsor, và được giải quyết bằng các tài liệu bổ sung ở Mục~\ref{sec:scrum-apply}.

\section{So sánh các phương pháp}

Bốn phương pháp được so sánh theo sáu tiêu chí rút ra từ các đặc điểm ở Bảng~\ref{tab:characteristics}
(Bảng~\ref{tab:methods}).

\begin{table}[htbp]
  \centering
  \caption{So sánh bốn phương pháp quản lý dự án}
  \label{tab:methods}
  \small
  \begin{tabularx}{\linewidth}{@{}P{3.3cm}LLLL@{}}
    \toprule
    \thead{Tiêu chí} & \thead{Waterfall} & \thead{Scrum} & \thead{Kanban} & \thead{Lai (Water-Scrum-Fall)} \\
    \midrule
    Yêu cầu thay đổi & Đắt, phải làm lại tài liệu & Rẻ, đưa vào sprint sau & Rẻ & Trung bình \\
    Thấy sản phẩm chạy được & Cuối dự án & Mỗi sprint & Liên tục & Giữa dự án \\
    Kiểm soát rủi ro kỹ thuật & Muộn & Sớm (POC Sprint 1, benchmark Sprint 2) & Sớm & Trung bình \\
    Kế hoạch, ngân sách cố định & Tốt & Theo velocity, cần bổ sung EVA & Yếu & Tốt \\
    Nhóm nhỏ, bán thời gian & Trung bình & Tốt (nhịp đều, mục tiêu ngắn) & Tốt, nhưng khó đo tiến độ theo mốc & Trung bình \\
    Đáp ứng báo cáo cho Sponsor & Tốt & Tốt khi bổ sung tài liệu quản lý & Yếu & Tốt \\
    \bottomrule
  \end{tabularx}
\end{table}

Theo Bảng~\ref{tab:methods}, Scrum và Kanban đều đón nhận thay đổi dễ dàng, kiểm soát rủi ro kỹ thuật sớm và hợp
với nhóm nhỏ làm bán thời gian. Scrum được chọn vì đo được tiến độ theo mốc và đáp ứng được yêu cầu báo cáo của
Sponsor khi bổ sung tài liệu quản lý, hai điểm Kanban còn yếu. Waterfall và mô hình lai mạnh về kế hoạch cố định
nhưng phát hiện rủi ro kỹ thuật muộn hơn.

\section{Cách áp dụng Scrum trong dự án}
\label{sec:scrum-apply}

\begin{itemize}
  \item \textbf{Vai trò:} Product Owner, Scrum Master kiêm Project Manager, và Developers (gồm cả BA và QC).
        Tổ chức nhóm và ma trận RACI xem~\cite{pmp}.
  \item \textbf{Sự kiện:} Sprint Planning (ước lượng bằng Planning Poker), Daily, Refinement, Sprint Review có
        demo và Retrospective.
  \item \textbf{Artefact:} Product Backlog~\cite{srs}; Sprint Backlog và burndown trên Azure Boards; Increment là
        nhánh feature chạy được, kèm test.
  \item \textbf{Tài liệu bổ sung cho Sponsor:} Charter, Business Case, WBS, CPM/Gantt, EVA và Risk Register là
        phần quản lý ``bên ngoài'' Scrum. Chúng được cập nhật mỗi sprint, không làm một lần rồi khoá.
  \item \textbf{Definition of Done:} áp dụng chung cho mọi story, định nghĩa trong Quality Plan~\cite{qtd}.
  \item \textbf{Quản lý mã nguồn:} mỗi nhóm việc một nhánh \texttt{feature/\ldots}, commit nhỏ theo task, tạo
        pull request trước khi merge vào \texttt{main}; thay đổi được ghi vào Change Log~\cite{mcr} và Release
        Notes~\cite{pcr}.
\end{itemize}

\section{Hạn chế và cách bù}

Scrum có bốn hạn chế đáng kể với một dự án phải báo cáo ngân sách và tiến độ cho Sponsor. Bảng~\ref{tab:limitations}
nêu cách dự án bù cho từng hạn chế.

\begin{table}[htbp]
  \centering
  \caption{Hạn chế của Scrum và biện pháp bù trong dự án}
  \label{tab:limitations}
  \begin{tabularx}{\linewidth}{@{}P{5.2cm}L@{}}
    \toprule
    \thead{Hạn chế của Scrum} & \thead{Cách bù trong dự án} \\
    \midrule
    Khó cam kết ngày kết thúc và ngân sách & Kết hợp velocity, lịch CPM~\cite{pmp} và EVA~\cite{mcr} \\
    Dễ phình phạm vi & Change Log có PO duyệt; phân quyền trong điều lệ (Bảng~\ref{tab:decision-rights}) \\
    Ít tài liệu & Bộ sáu tài liệu quản lý dự án; README và ARCHITECTURE trong repo \\
    Phụ thuộc kỷ luật họp & Daily ngắn; có thể họp không đồng bộ qua chat khi thành viên trùng lịch \\
    \bottomrule
  \end{tabularx}
\end{table}

% ============================================================ Chapter 5
\chapter{Kết luận và đề nghị phê duyệt}
\label{chap:conclusion}

\section{Tóm tắt}

Dự án SumTalk giải quyết việc ghi biên bản họp thủ công bằng một ứng dụng ghi âm, phiên âm, tóm tắt và dịch,
có thể tự host để giữ dữ liệu trong đơn vị. Hồ sơ này đã xác định:
\begin{itemize}
  \item sáu mục tiêu SMART, phạm vi và năm nhóm sản phẩm bàn giao, cùng sáu mốc gắn với các sprint;
  \item phương án tự xây (P1) là lựa chọn tốt nhất về tài chính: NPV khoảng +109,1 triệu VND, ROI 89,4\%, hoàn vốn
        sau khoảng 16 tháng, với điều kiện đơn vị có từ khoảng 23 cuộc họp mỗi tuần trở lên;
  \item Scrum là phương pháp quản lý, bổ sung bộ tài liệu quản lý để đáp ứng yêu cầu báo cáo của Sponsor.
\end{itemize}

\section{Việc cần làm trước khi trình duyệt}

\begin{enumerate}
  \item Điền các ô \canDien: người phụ trách, ngày phê duyệt, tổng giờ kế hoạch, độ dài sprint và ngày của các mốc.
  \item Thay giả định 600 giờ công bằng tổng giờ thực tế trên dashboard, rồi tính lại NPV, ROI và thời gian hoàn
        vốn bằng \texttt{pm\_calc.py}~\cite{pmcalc}.
  \item Lên kế hoạch đo số cuộc họp và thời gian tiết kiệm thật trong đợt UAT, vì đây là giả định có ảnh hưởng lớn
        nhất tới kết luận.
\end{enumerate}

\needspace{14\baselineskip}
\section{Phê duyệt}

Nhóm dự án đề nghị Sponsor phê duyệt khởi động dự án SumTalk theo phạm vi, ngân sách và phương pháp đã trình bày.

\begin{table}[H]
  \centering
  \renewcommand{\arraystretch}{2.4}
  \begin{tabularx}{\linewidth}{@{}P{3.6cm}LP{3.2cm}P{2.6cm}@{}}
    \toprule
    \thead{Vai trò} & \thead{Họ tên} & \thead{Chữ ký} & \thead{Ngày} \\
    \midrule
    Sponsor & & & \\
    Product Owner & & & \\
    Project Manager & & & \\
    \bottomrule
  \end{tabularx}
\end{table}

% ============================================================ References
\begin{thebibliography}{9}
\addcontentsline{toc}{chapter}{\bibname}
  \bibitem{srs} Nhóm SumTalk, \emph{Đặc tả yêu cầu phần mềm (Software Requirements Specification)},
    Tài liệu 02, mã SUMTALK-SRS, phiên bản 2.1, 24/11/2026.
  \bibitem{pmp} Nhóm SumTalk, \emph{Kế hoạch quản lý dự án (Project Management Plan)}, Tài liệu 03,
    mã SUMTALK-PMP, phiên bản 2.1, 24/11/2026.
  \bibitem{qtd} Nhóm SumTalk, \emph{Chất lượng và kiểm thử (Quality \& Test Document)}, Tài liệu 04,
    mã SUMTALK-QTD, phiên bản 2.1, 24/11/2026.
  \bibitem{mcr} Nhóm SumTalk, \emph{Báo cáo theo dõi và kiểm soát (Monitoring \& Control Report)}, Tài liệu 05,
    mã SUMTALK-MCR, phiên bản 2.1, 24/11/2026.
  \bibitem{pcr} Nhóm SumTalk, \emph{Báo cáo kết thúc dự án (Project Closure Report)}, Tài liệu 06,
    mã SUMTALK-PCR, phiên bản 2.1, 24/11/2026.
  \bibitem{pmcalc} Nhóm SumTalk, \emph{Công cụ tính NPV, ROI và thời gian hoàn vốn},
    kho mã nguồn \texttt{DatTapTanhCode/SumTalk}, tệp \path{docs/project/tools/pm_calc.py}.
\end{thebibliography}

\end{document}
